\documentclass[../main.tex]{subfiles}
\begin{document}
\markboth{MỞ ĐẦU}{MỞ ĐẦU}
\begin{center}
    {\large \textbf{MỞ ĐẦU}}
\end{center}
\vspace{0.4cm}

\noindent\textbf{Lý do chọn đề tài}

Trong bối cảnh xã hội hiện đại, khi nhịp sống ngày càng nhanh và áp lực công việc gia tăng, việc duy trì một chế độ ăn uống lành mạnh và cân bằng dinh dưỡng đang trở thành mối quan tâm lớn của nhiều người. Tuy nhiên, trên thực tế, phần lớn người dùng gặp khó khăn trong việc theo dõi và kiểm soát dinh dưỡng hàng ngày do thiếu thời gian, kiến thức chuyên môn cũng như công cụ hỗ trợ phù hợp. Việc ghi chép thủ công các bữa ăn hoặc ước lượng dinh dưỡng bằng cảm tính thường dẫn đến sai lệch, gây ảnh hưởng tiêu cực đến sức khỏe lâu dài.

Bên cạnh đó, sự phát triển của công nghệ trí tuệ nhân tạo (AI), đặc biệt trong lĩnh vực xử lý hình ảnh và nhận diện đối tượng, đã mở ra nhiều cơ hội mới trong việc hỗ trợ con người quản lý sức khỏe cá nhân. Việc ứng dụng AI để nhận diện món ăn từ hình ảnh giúp đơn giản hóa quá trình nhập liệu, giảm bớt gánh nặng cho người dùng và nâng cao độ chính xác trong việc phân tích giá trị dinh dưỡng.

Nhận thấy tiềm năng và tính thực tiễn cao của việc kết hợp công nghệ AI với quản lý dinh dưỡng, đề tài \textbf{``Xây dựng ứng dụng hỗ trợ chăm sóc sức khỏe và dinh dưỡng tích hợp AI''} được lựa chọn nhằm đáp ứng nhu cầu theo dõi chế độ ăn uống một cách khoa học, tiện lợi và hiệu quả, góp phần nâng cao chất lượng cuộc sống và ý thức chăm sóc sức khỏe trong cộng đồng.

\vspace{0.3cm}
\noindent\textbf{Mục tiêu của đồ án}

Đồ án hướng tới việc xây dựng một \textbf{ứng dụng quản lý dinh dưỡng thông minh}, hỗ trợ người dùng theo dõi, kiểm soát và điều chỉnh chế độ ăn uống hằng ngày một cách khoa học và hiệu quả. Hệ thống được phát triển nhằm tận dụng \textbf{công nghệ trí tuệ nhân tạo (AI)} trong việc nhận diện món ăn thông qua hình ảnh, từ đó tự động phân tích và ước lượng các chỉ số dinh dưỡng cần thiết cho cơ thể, giúp giảm thiểu các thao tác nhập liệu thủ công và nâng cao trải nghiệm người dùng.

Ứng dụng tích hợp các chức năng quan trọng như nhận diện món ăn từ hình ảnh, tính toán và hiển thị các thông tin dinh dưỡng cơ bản bao gồm năng lượng (calo), protein, chất béo, carbohydrate và một số vi chất liên quan. Bên cạnh đó, hệ thống hỗ trợ người dùng theo dõi tổng lượng dinh dưỡng nạp vào theo từng ngày và so sánh với các mục tiêu sức khỏe cá nhân như sử dụng các món ăn phù hợp với bệnh (tim, tiểu đường, mỡ máu, \dots), giảm cân, tăng cân hoặc duy trì thể trạng, từ đó giúp người dùng có cái nhìn tổng quan và chủ động điều chỉnh chế độ ăn uống phù hợp.

Ngoài các chức năng nghiệp vụ, đồ án cũng chú trọng đến việc \textbf{thiết kế giao diện thân thiện, dễ sử dụng}, phù hợp với nhiều đối tượng người dùng khác nhau. Đồng thời, việc nghiên cứu và áp dụng các công nghệ hiện đại trong phát triển ứng dụng không chỉ giúp nâng cao hiệu quả vận hành của hệ thống mà còn tạo tiền đề cho việc mở rộng và hoàn thiện trong tương lai.

Thông qua quá trình thực hiện đồ án, ứng dụng không chỉ góp phần nâng cao nhận thức của người dùng về dinh dưỡng và sức khỏe cá nhân, mà còn thể hiện tính ứng dụng thực tiễn của trí tuệ nhân tạo trong lĩnh vực chăm sóc sức khỏe, mang lại giá trị thiết thực và lâu dài cho người sử dụng.

\vspace{0.3cm}
\noindent\textbf{Bố cục của đồ án}

Nội dung của đồ án được trình bày trong các chương như sau:

\noindent\textbf{Chương 1: Cơ sở lý thuyết}
\begin{itemize}[label=$\checkmark$, leftmargin=1.2cm]
    \item Tổng quan về ứng dụng theo dõi và quản lý dinh dưỡng cá nhân
    \item Một số công nghệ khác
    \item Kết luận chương 1
\end{itemize}

\noindent\textbf{Chương 2: Xây dựng mô hình nhận diện món ăn}
\begin{itemize}[label=$\checkmark$, leftmargin=1.2cm]
    \item Giới thiệu nhiệm vụ
    \item Thu thập và tiền xử lý dữ liệu
    \item Lựa chọn mô hình
    \item Kết quả huấn luyện
    \item Đánh giá chi tiết
    \item Tích hợp mô hình
    \item Kết luận chương 2
\end{itemize}

\noindent\textbf{Chương 3: Phân tích, thiết kế hệ thống}
\begin{itemize}[label=$\checkmark$, leftmargin=1.2cm]
    \item Xác định yêu cầu
    \item Phân tích, thiết kế
    \item Thiết kế cơ sở dữ liệu
    \item Kết luận chương 3
\end{itemize}

\noindent\textbf{Chương 4: Cài đặt hệ thống}
\begin{itemize}[label=$\checkmark$, leftmargin=1.2cm]
    \item Yêu cầu hệ thống
    \item Một số kết quả cài đặt
    \item Kết luận chương 4
\end{itemize}

\vspace{0.3cm}
\noindent Các thành viên của nhóm thực hiện các nội dung đồ án chi tiết như sau:

\begin{enumerate}
    \item \textbf{Trần Bá Lợi -- B22DCCN510 (Nhóm trưởng)}
    \begin{itemize}
        \item Xây dựng đề cương, tài liệu đặc tả yêu cầu và theo dõi tiến độ chung của nhóm.
        \item Thu thập, gán nhãn và xây dựng cơ sở dữ liệu dinh dưỡng món ăn Việt Nam (tập VietFood68).
        \item Huấn luyện, đánh giá mô hình nhận diện món ăn (YOLOv11m) và xuất mô hình sang ONNX.
        \item Phân tích và thiết kế toàn bộ hệ thống.
        \item Xây dựng màn hình thống kê, biểu đồ và hồ sơ người dùng.
        \item Xây dựng các module: hồ sơ sức khỏe, nhật ký cân nặng/nước uống/vận động, thực đơn, thống kê tổng hợp hàng ngày.
        \item Xây dựng module lên thực đơn bằng AI.
        \item Kiểm thử tích hợp toàn hệ thống, sửa lỗi và hoàn thiện sản phẩm.
    \end{itemize}

    \item \textbf{Bùi Văn Nam - B22DCCN551}
    \begin{itemize}
        \item Phân tích và thiết kế toàn bộ hệ thống.
        \item Xây dựng RESTful API (Spring Boot): quản lý tài khoản, đăng ký/đăng nhập, xác thực JWT + Google OAuth2.
        \item Xây dựng hệ thống thông báo đẩy (FCM) và chức năng quét mã vạch tra cứu sản phẩm đóng gói.
        \item Xây dựng module quản lý món ăn.
        \item Tìm và thu thập các dữ liệu về món ăn trong hệ thống.
    \end{itemize}

    \item \textbf{Trần Đức Trung - B22DCCN875}
    \begin{itemize}
        \item Phân tích và thiết kế toàn bộ hệ thống.
        \item Xây dựng chatbot AI dựa theo hồ sơ người dùng.
        \item Xây dựng tính năng streak, huy hiệu, thử thách và kết bạn.
        \item Xây dựng các module về vận động của người dùng.
        \item Kiểm tra và rà soát toàn bộ hệ thống.
    \end{itemize}
\end{enumerate}

\end{document}
